\paragraph{Constants for the specified native trajectory.}
Write $n=N-1$. The limiting matrix of
$S_n^{-1}M(n)S_{n+1}$, where
$S_n=\operatorname{diag}(1,n^{-1},n^{-2},n^{-3})$, has an exact eigenvalue
$-1$, a positive eigenvalue $\Lambda$, and a nonreal conjugate pair of
modulus $\rho$. Exact polynomial isolation and rational logarithm enclosures give
\[
 \log\Lambda=53.28758332262020526735\ldots,\qquad
 \log\rho=-59.51995592483447507214\ldots .
\]
The accepted analytic rate for the selected zeta readouts is
\[
 E=\log\Lambda-\log\rho
   =112.80753924745468033949\ldots .
\]
The eigenvalue $-1$ is explicitly present; one must not identify $E$ with
the gap between the two largest moduli of the full matrix. The selected
readouts' elimination of that mode belongs to the analytic input.

For
\[
 C_n=\frac{(15n)!^2}{(10n)!^3},\qquad d_m=\operatorname{lcm}(1,\ldots,m),
\]
Stirling's formula gives $c_C=30\log(3/2)$. Hence the scaled coefficient
and error rates admit the safe enclosures
\[
 a=c_C+\log\Lambda<65.451536565866,\qquad
 b=c_C+\log\rho<-47.356002681589.
\]
The alternative factor $(14n)!(16n)!/(10n)!^3$ has exponent
$14\log14+16\log16-30\log10=12.23066938062874988548\ldots$,
instead of $12.16395324324493145934\ldots$ for $C_n$.
Any prime-divisor profile must be changed consistently when switching
between these normalizations.

\paragraph{The two arithmetic regimes.}
Put
\[
 q_n=W_{00},\quad
 p_{3,n}=-\frac12W_{20}+\frac14W_{30},\quad
 p_{4,n}=-\frac16W_{30}.
\]
These are exactly the user's numerator and denominator divided by $12$.
Suppose the arithmetic argument supplies divisors $G_{3,n}$ and $G_{4,n}$
with logarithmic densities at least $\gamma_3,\gamma_4$, respectively, and
\[
 \frac{6C_nd_{19n}^3}{G_{3,n}}(p_{3,n},q_n)\in\mathbb Z^2,
 \qquad
 \frac{6C_nd_{19n}^4}{G_{4,n}}(p_{3,n},p_{4,n},q_n)\in\mathbb Z^3.
\]
Fixed integer coefficients $A,B$ introduce no new exponential factor.
For $B=0$ one uses the first normalization. For $B\ne0$, one uses the
second and forms $Ap_{3,n}+Bp_{4,n}$; the triangle inequality preserves
the accepted error upper bound, and cancellation can only improve it.
Fixed rational coefficients are handled by one fixed denominator multiplier.
The resulting safe growth and decay exponents are
\[
\begin{array}{c|cc}
 &\sigma\text{ upper bound}&\tau\text{ lower bound}\\\hline
 B=0&122.451536565866-\gamma_3&\gamma_3-9.643997318411\\
 B\ne0&141.451536565866-\gamma_4&\gamma_4-28.643997318411.
\end{array}
\]
Whenever the respective $\tau$ is positive, the accepted analytic
error and growth inputs give $\liminf\delta_N\ge\tau/\sigma$.
This implication requires a proved uniform divisor density. A finite gcd
or fitted denominator slope cannot be substituted for $\gamma_3$ or
$\gamma_4$. No irrationality-measure upper bound is asserted here; that
would require suitable additional nonvanishing and spacing hypotheses.

\paragraph{An elementary baseline behind the divisor improvement.}
The normalized rational kernel is a product of three binomial polynomials
of degree $10n$ and two reciprocal blocks
$m!/(t+u)_{m+1}$ with $m=15n$.
For a binomial polynomial of degree $m$, its Taylor coefficient of degree
$j$ at an integer becomes integral after multiplication by $d_m^j$.
Indeed, Vandermonde's identity reduces this statement to
$\binom{x}{l}=(-1)^{l-1}(x/l)\prod_{i=1}^{l-1}(1-x/i)$.
A reciprocal block has integer residues, by its explicit binomial
partial fractions. The difference between any two pole indices in this
case is at most $16n$. Therefore the double-pole coefficients $a_j$ and
simple-pole coefficients $b_j$ satisfy
\[
 C_na_j\in\mathbb Z,\qquad d_{16n}C_nb_j\in\mathbb Z.
\]
The finite readouts involve $H_{j+2n}^{(3)},H_{j+2n}^{(2)}$ for $\zeta(3)$,
and $H_{j+2n}^{(4)},H_{j+2n}^{(3)}$ for $\zeta(4)$, with $j+2n\le19n$.
This yields the displayed normalizations with $G_{3,n}=G_{4,n}=1$.
The extra prime divisors are essential for a positive mixed decay exponent;
the baseline alone does not establish the desired irrationality bound.
